\documentclass[12pt]{article}
\usepackage[left=2cm,right=2cm,top=2cm,bottom=2cm]{geometry}
\usepackage{wrapfig}
\usepackage{booktabs}
\usepackage{graphicx}
\usepackage{amssymb}
\usepackage{siunitx} % Required for alignment
\usepackage{subfigure}
\usepackage{multirow}
\usepackage{rotating}
\usepackage[T2A]{fontenc}
\usepackage[russian]{babel}
\usepackage{caption}
\usepackage{hyperref}
\usepackage{mathtools}
\usepackage{amsmath}
\usepackage{float}

\hypersetup{
    colorlinks=true,
    linkcolor=blue,
    urlcolor=magenta,
    citecolor=blue,
    }
\graphicspath{{pictures/}}

\begin{document}
  \begin{titlepage}
      \begin{center}
          \vspace*{1cm}

          \Huge
          \textbf{ДИАМАГНЕТИЗМ И ПАРАМАГНЕТИЗМ}

          \vspace{1.5cm}

          \Large
          \textbf{Балдин Виктор Б01-303}

          \vfill

          Вопрос по выбору \\
          Устный экзамен по общей физике

          \vspace{0.8cm}

          \includegraphics[width=0.4\textwidth]{university_logo.png}

          Физтех-школа радиотехники и компьютерных технологий\\
          Московский физико-технический институт\\
          Долгопрудный, 2024
      \end{center}
  \end{titlepage}

  \tableofcontents

  \section{Введение}

  \section{Диамагнетизм}

  \section{Парамагнетизм}

  \section{Заключение}

% \bibliography{references}
% \bibliographystyle{gost}

\end{document}
